\section[Dilemme de l'infrastructure]{Le dilemme de l'infrastructure: entre le sur-mesure et le standard}
La refonte de la bibliothèque numérique de la BIU Cujas soulève une interrogation stratégique dépassant la simple technique informatique : quel modèle d'infrastructure privilégier pour garantir à la fois l'indépendance et l'adéquation aux besoins scientifiques et la pérennité de l'outil ? En ingénierie documentaire, c'est une question classique opposant la logique du développement spécifique, incarné par le prototype Codex, à celle de l'adoption d'outil standard, comme Omeka S.

\subsection{L'analyse du sur-mesure}
Le choix initial de la bibliothèque Cujas, en développant le prototype Codex, offre une flexibilité totale et une forte adaptation aux besoins métiers. Cela évite également la dépendance finiancière et technique vis-à-vis d'un prestataire externe, \enquote{qui facture cher la moindre évolution}. Conçu en Python-Django, Codex offre l'avantage de s'interfacer dans la façon de travailler des experts métiers. Ce couplage permettrait  d'automatiser certaines tâches de la chaîne de numérisation de Cujas, comme la génération de métadonnées METS ou XML, ou des transferts sécurisés vers le serveur d'archivage du CINES.

Si l'adaptation aux besoins est forte, ce mode de développement démontre une extrême vulnérabilité humaine, identifiée, par le directeur de la bibliothèque. Il le perçoit comme un risque de pérennité de l'infrastructure en bâtissant cette future bibliothèque numérique uniquement sur les compétences d'une seule personne. En effet, Olesea Dubois confirme cette analyse : \enquote{le plus gros soucis avec les outils maisons c'est qu'ils sont très souvent associés àdes personnes et à des procédures qui leur sont prores}. Un autre inconvénient mis en avant par la responsable de la bibliothèque numérique de Sciences Po est que \enquote{généralement, il y a peu ou pas assez de documentations entraînant ainsi une durée de vie assez courte de la bibliothèque numérique}. Effectivement, le peu de documentation disponible entraîne un risque pour la pérennité de l'outil, car sans documentation cela empêche le maintien sur le long terme de la bibliothèque numérique. avoir recours à un logiciel sur-mesure entraîne une difficulté de recrutement, car l'institution se coupe d'un vivier de compétence, rendant rapidement le code sur-mesure inexploitable en cas du départ de son concepteur condamnant le système à une obsolescence rapide. La crainte de la plupart des institutions patrimoniales y compris la bibliothèque Cujas est de reproduire la situation qui a entraîné la disparition de la bibliothèque numérique Medica. C'est pour cela qu'elles ont choisi de se tourner vers des outils standardisés afin de disposer d'un vivier de personnel en cas de départ de l'un des concepteurs de la bibliothèque numérique.

\subsection{L'analyse du standard communautaire}
 Face au risque d'isolement technique du sur-mesure, les retours d'expérience des responsables de bibliothèque numérique des institutions patrimoniales comme la Bibliothèque Sainte-Geneviève, la Bibliothèque de Sciences Po ou l'ENS, démontre un abandon progressif des outils sur-mesure au profit d'Omeka S. Financé par des subventions de recherche, de fondations et soutenu par une communauté, Omeka S offre des garanties majeures de durabilité. Contrairement à Django qui nécessitant d'avoir des compétences informatiques avancées pour développer le logiciel et l'intégrer dans le web sémantique. La particularité d'Omeka S est d'avoir été pensé pour s'intégrer directerement ou plus facilement dans le web de données. Il s'appuie sur une \enquote{modélisation des métadonnées sous RDF} combinéeés à des vocabulaires standardisées, comme le Dublin Core, favorise \enquote{une meilleure exposition des données}. Il facilite l'exposition des métadonnées, leur exportation et leur interopérabilité. 
 
 De plus, un autre des avantages pour les catalogueurs est de \enquote{bénéficier d'un back-office documenté, collaboratif et modulaire}. Toutefois, malgré les avantages certains, l'environnement Omeka S s'accompagne de contraintes réelles. Bien qu'open-source, le logiciel est \enquote{gratuit}, mais il exige des ressources pour fonctionner et évoluer. Omeka S demande des ressources internes pour en assurer l'hébergement, le paramétrage et la maintenance technique. Gilles Clément, l'un des concepteurs de NuBIS, témoigne de cette limite : \enquote{À l'usage, chaque nouvelle version d'Omeka S cassait quelque chose...Nous étions confrontés à beaucoup de bidouillages en raison de l'instabilité des modules tiers développés par des prestataires différents, ce qui rendait le code lourd et difficile à maintenir}.
 
 Face à cet obstacle, la solution mise en place par la BIS pour le résoudre, a été de séparer techniquement le front du back, \enquote{rendant ainsi le code-source plus maintenable dans le temps}. En effet, à chaque mise à jour, le module complémentaire déréglait les paramétrages du front-office. la BIS a donc découplé son architecture en décidant de \enquote{développer leur back-office sous Omeka S et du front sous Astro.js}. Cette distinction rend étanche les couches logicielles permettant ainsi le bon fonctionnement de NuBIS et la satisfaction des équipes qui en sont en charge. Malgré cet obstacle rencontré par certaines institutions, l'open-source, contrairement au sur-mesure, a cette \enquote{notion de mutualisation rendant plus facile l'idée du maintien de cette structure}. Le choix de la BIS du framework Astro.js répond à des exigences d'éco-conception et d'accessibilité numériqe, évitant ainsi l'interrogation de la base à chaque clique de l'usager. L'avantage majeur de ce type de fonctionnement est que les montées de versions d'Omeka S n'impactent plus l'affichage public, sécurisant mieux le back-office le rendant plus maintenable dans le temsp et pérenne.